\subsection{Regime 4: saturating antibody}

In this regime, antibody is so abundant that even the weak, off-target
interactions are saturated. As in Regime 1, let $S\in\{S_c,S_t\}$ with
$S_c=A+B$ and $S_t=A'+B$, and assume the following for both conditions.

\textbf{Assumption.} The antibody left over after binding \emph{every}
epitope, target and off-target, is much greater than $K_{d2}$:
$$
Y-S-(C+D)\gg K_{d2}.
$$
For brevity, write $E=Y-S-(C+D)$ for this \textbf{excess antibody}
($E_c$ and $E_t$ for the two conditions). Because $S_t<S_c$ when $A'<A$,
$$
E_t=Y-S_t-(C+D)>Y-S_c-(C+D)=E_c,
$$
so it suffices to check the assumption for the control condition,
$E_c\gg K_{d2}$.

\textbf{Intuition.} Each epitope can hold at most one antibody, so at
most $S+(C+D)$ antibody can be bound, and at least $E$ must remain free.
Recall that $K_{d2}$ is the free-antibody concentration at which
off-target epitopes are half-bound. If the guaranteed free antibody $E$
is already far above $K_{d2}$, then off-target epitopes are nearly all
bound---and target epitopes, which bind more tightly ($K_{d1}<K_{d2}$),
are even more completely bound. Note that this requires both $Y\gg K_{d2}$
\emph{and} $Y>S+C+D$: antibody must outnumber all epitopes, including the
typically very abundant off-target ones.

We now derive $\rho$ for either condition. (For the treatment condition,
$Y_f$, $\theta_1$, $\theta_2$, $E$, and $\rho$ should be read as their
primed counterparts.)

\textbf{Step 1: a lower bound on $Y_f$.} Start from antibody conservation
(derived in Regime 1),
$$
Y=Y_f+S\theta_1+(C+D)\theta_2.
$$
Solving for $Y_f$, and then adding and subtracting $S+(C+D)$,
$$\begin{aligned}
Y_f
&=Y-S\theta_1-(C+D)\theta_2 \\
&=\bigl[Y-S-(C+D)\bigr]+S-S\theta_1+(C+D)-(C+D)\theta_2 \\
&=E+S(1-\theta_1)+(C+D)(1-\theta_2).
\end{aligned}$$
Intuitively, the free antibody is the excess antibody $E$, plus one
antibody ``returned to solution'' for each epitope that is unbound.
Because $\theta_1\le1$ and $\theta_2\le1$, the last two terms are
nonnegative, so
\begin{equation}
Y_f\ge E.
\label{eq:r4-1}
\end{equation}
By the assumption, $E\gg K_{d2}$, so $Y_f\gg K_{d2}$. Since $K_{d1}<K_{d2}$,
also $Y_f\gg K_{d1}$.

\textbf{Step 2: both classes of epitope are nearly fully bound.} From
$\theta_2=Y_f/(Y_f+K_{d2})$,
$$
1-\theta_2
=\frac{Y_f+K_{d2}}{Y_f+K_{d2}}-\frac{Y_f}{Y_f+K_{d2}}
=\frac{K_{d2}}{Y_f+K_{d2}}.
$$
Dropping $K_{d2}$ from the denominator makes the fraction larger, and
then substituting \eqref{eq:r4-1} (a larger denominator makes it smaller),
$$
1-\theta_2
\le\frac{K_{d2}}{Y_f}
\le\frac{K_{d2}}{E}
\ll1.
$$
By the same logic, $1-\theta_1=K_{d1}/(Y_f+K_{d1})\le K_{d1}/E$, which is
even smaller because $K_{d1}<K_{d2}$. Thus $\theta_1\approx\theta_2\approx1$:
nearly every epitope, target or off-target, is bound by antibody.

\textbf{Step 3: an exact expression for $1-\rho$.} Writing $1$ with the
same denominator as $\rho$,
\begin{equation}
1-\rho
=\frac{Y_f+K_{d2}}{Y_f+K_{d2}}-\frac{Y_f+K_{d1}}{Y_f+K_{d2}}
=\frac{K_{d2}-K_{d1}}{Y_f+K_{d2}}.
\label{eq:r4-2}
\end{equation}
Since $K_{d2}>K_{d1}$, the numerator is positive, so $1-\rho>0$, i.e.,
$\rho<1$.

\textbf{Step 4: $\rho\approx1$.} In \eqref{eq:r4-2}, the numerator
satisfies $K_{d2}-K_{d1}<K_{d2}$, and the denominator satisfies
$Y_f+K_{d2}>Y_f\ge E$ by \eqref{eq:r4-1}. A larger numerator and a
smaller denominator each make a positive fraction larger, so
\begin{equation}
0<1-\rho<\frac{K_{d2}}{Y_f}\le\frac{K_{d2}}{E}=\frac{K_{d2}}{Y-S-(C+D)}\ll1.
\label{eq:r4-3}
\end{equation}
Therefore $\rho\approx1$, and by the same argument for the treatment
condition, $\rho'\approx1$.

\textbf{Step 5: comparing $\rho'$ and $\rho$.} First, $\rho$ increases
with $Y_f$. Differentiating with the quotient rule,
$$
\frac{d\rho}{dY_f}
=\frac{(1)(Y_f+K_{d2})-(Y_f+K_{d1})(1)}{(Y_f+K_{d2})^2}
=\frac{K_{d2}-K_{d1}}{(Y_f+K_{d2})^2}
>0.
$$
Second, $Y_f'>Y_f$ when $S_t<S_c$. To see this, substitute
$\theta_1=Y_f/(Y_f+K_{d1})$ and $\theta_2=Y_f/(Y_f+K_{d2})$ into antibody
conservation, and define the right-hand side as a function of a trial
free-antibody concentration $x\ge 0$ and the target total $S$:
$$
F(x,S)=x+S\frac{x}{x+K_{d1}}+(C+D)\frac{x}{x+K_{d2}}.
$$
The equilibrium $Y_f$ is the value of $x$ that solves $F(x,S)=Y$.
\begin{itemize}
    \item $F$ is increasing in $x$: each term is increasing, since
    $\frac{d}{dx}x=1>0$ and, by the quotient rule, for any $K>0$,
    $$
    \frac{d}{dx}\frac{x}{x+K}
    =\frac{(1)(x+K)-x(1)}{(x+K)^2}
    =\frac{K}{(x+K)^2}>0.
    $$
    \item $F$ is increasing in $S$ for $x>0$: $\partial F/\partial S
    =x/(x+K_{d1})=\theta_1>0$.
\end{itemize}
The two conditions share $Y$, so $F(Y_f,S_c)=Y=F(Y_f',S_t)$. Because
$S_t<S_c$ and $Y_f>0$ (by \eqref{eq:r4-1}), the second property gives
$$
F(Y_f,S_t)<F(Y_f,S_c)=Y=F(Y_f',S_t).
$$
Since $F(\cdot,S_t)$ is increasing, $F(Y_f,S_t)<F(Y_f',S_t)$ requires
$Y_f<Y_f'$. In words: with fewer target epitopes to absorb antibody in
the treatment condition, more antibody is left free.

Combining the two facts, $Y_f'>Y_f$ implies $\rho'>\rho$, just as in
Regime 1. But both lie below 1 (Step 3), so using \eqref{eq:r4-3},
$$
0<\rho'-\rho<1-\rho<\frac{K_{d2}}{E_c}\ll1.
$$
The difference is negligible: $\rho\approx\rho'\approx1$.

\textbf{Result.} The immunoprecipitation contains
$A\theta_1:B\theta_1:C\theta_2:D\theta_2=A:B:\rho C:\rho D$, so
$$
\text{IP composition}\approx A:B:C:D,
$$
which is the same as the input. Because nearly every epitope is captured
(Step 2), the immunoprecipitation is essentially a copy of the
epitope-containing chromatin, with no preference for target over
off-target epitopes.

\subsubsection{ChIP normalization}

We reuse the read model and normalization formulas from Regime 1,
substituting $\rho\approx\rho'\approx1$.

\textbf{Scaling by spike-in reads (ChIP-Rx).} From Regime 1, the ratio
between conditions is
$$
\frac{n_i'/N_s'}{n_i/N_s}
=\frac{a_i'}{a_i}\cdot\frac{B+\rho D}{B+\rho' D}.
$$
Since $\rho'>\rho$ (Step 5), $B+\rho'D>B+\rho D$, so the bias factor is
slightly less than 1. To bound how far below 1 it can be, use
$\rho'<1$ to enlarge the denominator, $B+\rho'D<B+D$, which makes the
fraction smaller:
$$\begin{aligned}
\frac{B+\rho D}{B+\rho' D}
&>\frac{B+\rho D}{B+D}
=\frac{(B+D)-(1-\rho)D}{B+D}
=1-(1-\rho)\frac{D}{B+D}.
\end{aligned}$$
Since $D/(B+D)\le1$ and, by \eqref{eq:r4-3}, $1-\rho<K_{d2}/E_c$,
$$
1-\frac{K_{d2}}{E_c}<\frac{B+\rho D}{B+\rho' D}<1.
$$
So the bias factor is $\approx1$, and \textbf{ChIP-Rx normalization is
essentially unbiased} at target loci in this idealized model: both
conditions capture both spike-in epitope classes with the same (nearly
complete) efficiency, so the spike-in reads track the amount of chromatin
just as they would in the input.

\textbf{Scaling by total reads (RPM).} Substituting $\rho\approx\rho'\approx1$
into the Regime 1 ratio,
$$
\frac{n_i'/N'}{n_i/N}
=\frac{a_i'}{a_i}\cdot\frac{A+B+\rho(C+D)}{A'+B+\rho'(C+D)}
\approx\frac{a_i'}{a_i}\cdot\frac{A+B+C+D}{A'+B+C+D}.
$$
If the off-target pool dominates ($C+D\gg A$), the second factor is
close to 1 and even RPM recovers the global change; the large,
unchanging off-target pool acts as an internal reference. This is not a
virtue, as we now see.

\textbf{The cost: loss of specificity.} Let $c_j$ be the off-target
epitope concentration at an off-target locus $j$. The immunoprecipitation
contains $a_i\theta_1$ at target locus $i$ and $c_j\theta_2$ at locus
$j$. Dividing both by $\theta_1$ and substituting $\rho=\theta_2/\theta_1$,
$$
\frac{a_i\theta_1}{c_j\theta_2}
=\frac{a_i}{\rho\,c_j}
\approx\frac{a_i}{c_j},
$$
which is the same ratio as in the input. Compare Regime 1, where
$\rho\approx r$ boosts target over off-target by a factor $1/r=K_{d2}/K_{d1}\gg1$.
Here that enrichment is lost entirely. Summing over loci, the fraction
of IP reads that come from endogenous off-target epitopes is
$$
\frac{\rho C}{A+B+\rho(C+D)}\approx\frac{C}{A+B+C+D},
$$
and the fraction from all off-target epitopes is approximately
$(C+D)/(A+B+C+D)$. When off-target epitopes are far more abundant than
the target (as is typical, e.g., for an antibody that weakly recognizes
unmodified histones), most reads are off-target and peaks at target loci
are diluted into the background. Normalization is unbiased, but the
signal it normalizes is mostly noise.

\textbf{In practice.} This regime is an idealization. Real ChIP
``background'' includes not only weak antibody-epitope binding (which
saturates, as modeled here) but also chromatin that sticks
nonspecifically to beads, tubes, or antibody-bead complexes. That
background does not follow the binding isotherm $Y_f/(Y_f+K_{d2})$; it
roughly scales with the amount of input chromatin and need not be the
same in the endogenous and spike-in chromatin. Moreover, reaching
$Y>S+C+D+K_{d2}$ requires enormous amounts of antibody when the
off-target pool is large. The main lesson of this regime is therefore
not that saturating antibody makes spike-in normalization safe, but that
adding antibody beyond the point of saturating the target trades
specificity for (at best) uniform capture: the immunoprecipitation
increasingly resembles the input.
